\section{Conclusion}
\label{sec:conclusion}

We have presented a reinforcement-learning framework for training small
language models to deliberate on verifiable problems, in which the central
learnable object is a parametric action-value function $\Qpsi(s, a)$ over
an LLM-induced context MDP, MCTS serves as a policy improvement operator
that supplies stronger Bellman targets $\Qmcts(s, a)$ during training and
as an optional test-time-compute knob at inference, and reward enters only
through verifier-anchored signals so that the Goodhart attack surface
associated with learned-aesthetic critics is closed by construction. The
six contributions are an approximate-$Q$-learning framing for the
verifiable context MDP, a shared-encoder architecture with a frozen base
model and trainable multi-head $Q$ and policy heads, a reward principle
that quarantines learned signals away from Bellman backups, an explicit
role for MCTS as a sample-efficient policy improvement operator and a
test-time-compute knob rather than as the central learning algorithm,
a pre-registered experimental plan whose baseline grid and headline
frontier figure are fixed before any training begins, and a preliminary
feasibility probe that exercised the framework's integration end-to-end
before the benchmark study begins.

The probe (Section~\ref{sec:pilot}) established feasibility rather than a
result. It predates the current action space, so we report no accuracy or
agreement number from it; what it contributes is the demonstration that every
component of the framework can be implemented and the expert-iteration loop run
end-to-end, together with the methodological precautions the pre-registered
study inherits as monitoring requirements. Those precautions are concrete:
train/held-out problems must share structural shape but never surface content;
per-round training-set composition and per-class held-out accuracy must be
reported alongside any multi-round result; target stability must be measured as
pairwise recommended-verb agreement across independent searches from a fixed
state, stratified by problem family, with the heavy-search budget set by the
slowest-converging family; and the compute axis must charge every LLM forward
inside the search, with the run harness persisting the separated learning
signals per iteration.

By pre-registering the baseline grid, the compute axis, and the
falsification criteria (Section~\ref{sec:experimental-plan}), we
constrain the eventual results section to test the framework rather than
to fit to it: a frontier that fails to clear the best-of-$N$ baseline at
matched compute is reported as such, not retrofitted into a different
claim. The pre-registered falsifiable hypothesis is precisely that the
heavy-search B4 curve lies strictly above the best-of-$N$ B1 curve at
matched inference compute on the benchmarks where both produce signal
(Section~\ref{sec:exp-frontier}); the larger reasoning model B5 is a
reference point, not a target, and the distance between B4 and B5
calibrates how much of the headline gap the loop has closed rather than
serving as a bar the framework must clear. Whether B4 clears B1, and how
far it narrows the B4-to-B5 distance, is the question the implementation
will answer. The pilot has shown that the framework's mechanics are
not vacuous; whether they generalize to standardized benchmarks under the
pre-registered protocol is the next empirical question.
